\section{实验结果与分析}

\subsection{实验设置}

\subsubsection{实验数据集}

\subsubsection{实验评估指标}

\subsubsection{实现细节}

\subsection{实验结果分析}


\subsection{消融实验与讨论}

以下是 \LaTeX 和 Microsoft Word 的优劣比较：

\begin{table}[h] % 使用 [h] 强制表格位置
  \centering
    \bicaption{\LaTeX  与 Word 优劣比较}{This is the English title}
  \label{tab:latex-vs-word} % 表格标签
  \begin{tabular}{lcc} % 定义表格列格式（3 列，第一列左对齐，其余居中）
    \toprule % 顶部横线
    \textbf{比较项} & \textbf{\LaTeX} & \textbf{Word} \\
    \midrule % 中间横线
    排版质量 & 高 & 中 \\
    易用性 & 低 & 高 \\
    功能与灵活性 & 高 & 中 \\
    协作与兼容性 & 中 & 高 \\
    适用场景 & 学术论文、技术文档 & 日常办公、简单文档 \\
    \bottomrule % 底部横线
    
  \end{tabular}
\end{table}
      \tablesource{这是表注}

如表 \ref{tab:latex-vs-word} 所示，LaTeX 和 Word 在排版质量、易用性、功能与灵活性、协作与兼容性以及适用场景等方面各有优劣\cite{bai2002kuo}\cite{xu2024lun}。

\subsection{本章小结}